\documentclass[11pt, a4paper]{article}

\usepackage[T1]{fontenc}
\usepackage[utf8]{inputenc}
\usepackage{tgpagella}
\usepackage[spanish, es-noshorthands]{babel}
\usepackage{amsmath, amssymb}
\usepackage{booktabs}
\usepackage{tabularx}
\usepackage{multicol}
\usepackage[most]{tcolorbox}
\usepackage[american]{circuitikz}
\usepackage[margin=2.5cm]{geometry}
\usepackage{fancyhdr}

\pagestyle{fancy}
\fancyhf{}
\setlength{\headheight}{16pt}
\lhead{\textbf{UCB Sede Tarija} | Ing. Mecatrónica}
\chead{Circuitos Electrónicos II (IMT-221)}
\rhead{Recuperatorio 1 (EC1)}
\lfoot{Evaluación Sumativa}
\rfoot{Página \thepage}
\renewcommand{\headrulewidth}{0.4pt}
\definecolor{ucbblue}{RGB}{0, 51, 102}

\begin{document}

\begin{center}
    \Large\textbf{\textcolor{ucbblue}{Recuperatorio 1: Análisis Dinámico de CA}}[cite: 9]
\end{center}

\begin{tcolorbox}[colback=ucbblue!5!white, colframe=ucbblue, title=\textbf{Instrucciones y Condiciones de Evaluación[cite: 9]}]
    \begin{itemize}
        \setlength{\itemsep}{0pt}
        \item \textbf{Tiempo de aplicación:} 80 minutos[cite: 9].
        \item \textbf{Material permitido:} Calculadora científica y formulario de una hoja[cite: 9].
        \item \textbf{Convención:} Todos los fasores de tensión están expresados en valor pico ($\hat{V}$)[cite: 9].
        \item Muestre todo el desarrollo analítico. Un resultado final sin procedimiento no es válido[cite: 9].
        \item Conserve las cantidades complejas durante todo el análisis y exprese los resultados solicitados en magnitud y fase[cite: 9].
    \end{itemize}
\end{tcolorbox}

\vspace{0.4cm}
\section*{Problema 1: Red RLC Cargada (25 puntos)[cite: 9]}
Una fuente sinusoidal ideal $\hat{V}_s = 10\angle0^\circ\,\text{V}_{\text{pico}}$ opera a $f = 1.00\,\text{kHz}$[cite: 9]. La fuente alimenta primero un capacitor $C = 159.15\,\text{nF}$[cite: 9]. Después del capacitor se encuentra el nodo $A$[cite: 9]. Entre el nodo $A$ y tierra existen dos ramas en paralelo:
\begin{itemize}
    \setlength{\itemsep}{0pt}
    \item \textbf{Rama 1:} Resistor $R_1 = 1.00\,\text{k}\Omega$[cite: 9].
    \item \textbf{Rama 2:} Resistor $R_2 = 1.00\,\text{k}\Omega$ en serie con un inductor $L = 159.15\,\text{mH}$[cite: 9].
\end{itemize}
El nodo $B$ se encuentra entre $R_2$ y $L$[cite: 9]. La salida solicitada es $\hat{V}_B$, la tensión sobre el inductor referida a tierra[cite: 9].

\begin{center}
    \begin{circuitikz}[scale=1.2, transform shape, thick]
        \draw (0,0) to[V, v=$\hat{V}_s$] (0,3) to[C, l=$C$] (3,3) coordinate(A);
        \draw (A) to[R, l_=$R_1$] (3,0);
        \draw (A) -- (5,3) to[R, l=$R_2$] (5,1.5) coordinate(B);
        \draw (B) to[L, l=$L$, v=$\hat{V}_B$] (5,0);
        \draw (0,0) -- (5,0) node[ground]{};
        
        \node[above] at (A) {$A$};
        \node[right] at (B) {$B$};
    \end{circuitikz}
\end{center}

\textbf{a) Modelo fasorial (6 pts):} Determine $Z_C$ y $Z_L$ a la frecuencia indicada[cite: 9]. Identifique la impedancia total de la rama $R_2 + L$[cite: 9]. \\
\textbf{b) Modelo equivalente y método (6 pts):} Determine la impedancia equivalente vista desde el nodo $A$ ($Z_A$) y la impedancia total vista por la fuente ($Z_T$)[cite: 9]. Indique brevemente por qué su método es apropiado[cite: 9]. \\
\textbf{c) Análisis fasorial (9 pts):} Determine $\hat{V}_B$ en forma rectangular y polar, mostrando los pasos intermedios[cite: 9]. \\
\textbf{d) Interpretación (4 pts):} Respecto de $\hat{V}_s$, ¿la magnitud de $\hat{V}_B$ es mayor o menor?[cite: 9] ¿$\hat{V}_B$ adelanta o atrasa a la fuente?[cite: 9] Explique qué representa físicamente ese ángulo[cite: 9].

\newpage
\section*{Problema 2: Respuesta Frecuencial RLC (25 puntos)[cite: 9]}
Una fuente sinusoidal alimenta una red serie formada por $R = 1.00\,\text{k}\Omega$, $L = 159.15\,\text{mH}$ y $C = 159.15\,\text{nF}$[cite: 9]. La salida $\hat{V}_o$ se mide únicamente sobre $R$[cite: 9].

\begin{center}
    \begin{circuitikz}[scale=1.2, transform shape, thick]
        \draw (0,0) to[V, v=$\hat{V}_s$] (0,2) to[L, l=$L$] (2,2) to[C, l=$C$] (4,2) to[R, l=$R$, v=$\hat{V}_o$] (4,0) -- (0,0);
    \end{circuitikz}
\end{center}

Defina la función de transferencia como $H(j\omega) = \frac{\hat{V}_o}{\hat{V}_s}$[cite: 9].

\vspace{0.3cm}
\textbf{a) Función de transferencia (8 pts):} Derivando desde las impedancias, encuentre $H(j\omega)$[cite: 9]. Simplifique su respuesta hasta obtener una expresión en términos de $R, L, C$ y $\omega$[cite: 9]. \\
\textbf{b) Frecuencia característica (5 pts):} Determine la frecuencia $f_0$ para la cual las reactancias inductiva y capacitiva se cancelan[cite: 9]. Evalúe $H(j\omega_0)$ e interprete físicamente el resultado[cite: 9]. \\
\textbf{c) Evaluación fuera de $f_0$ (8 pts):} Evalúe la función a $f = 500\,\text{Hz}$[cite: 9]. Determine $H(j\omega)$ en forma rectangular, su magnitud $|H|$, su fase $\angle H$ y la ganancia en decibelios $20\log_{10}|H|$[cite: 9]. \\
\textbf{d) Interpretación frecuencial (4 pts):} Sin recalcular, explique cualitativamente qué espera para frecuencias $f \ll f_0$ y $f \gg f_0$[cite: 9]. Clasifique cualitativamente la respuesta del circuito (ej. pasabajas, pasaaltas, etc.)[cite: 9].

\vspace{2cm}
\begin{tcolorbox}[colback=white, colframe=black, title=\textbf{Auditoría Autónoma (Antes de entregar)[cite: 9]}]
    \begin{multicols}{2}
    \begin{itemize}
        \setlength{\itemsep}{0pt}
        \item[$\square$] ¿Modelé correctamente $L$ y $C$?[cite: 9]
        \item[$\square$] ¿Interpreté correctamente serie/paralelo?[cite: 9]
        \item[$\square$] ¿Conserve las polaridades definidas?[cite: 9]
        \item[$\square$] ¿Mis ecuaciones provienen del circuito?[cite: 9]
        \item[$\square$] ¿Operé cantidades complejas debidamente?[cite: 9]
        \item[$\square$] ¿Mis unidades son correctas?[cite: 9]
        \item[$\square$] ¿Respondí en magnitud y fase?[cite: 9]
        \item[$\square$] ¿La fase es coherente con el circuito?[cite: 9]
    \end{itemize}
    \end{multicols}
\end{tcolorbox}

\end{document}
